%% RNAAS research note: new DBV WD J203728.70+145353.15
%% Compile with AASTeX 6.3.1 (e.g. on Overleaf: upload this file, fig1.pdf and note.bib).
%% Numbers: data/labs/dbv_note/fits.json, data/labs/dbv_ztf/ztf_result*.json (wdplanets/labs/dbv_note.py, dbv_ztf.py).
\documentclass[RNAAS]{aastex631}

\begin{document}

\title{WD J203728.70+145353.15: a New Pulsating DB White Dwarf at 88 pc Found in TESS and Confirmed with ZTF}

\author{Christian Jair}
\affiliation{Independent researcher}

\begin{abstract}
We report pulsations in the DB white dwarf WD J203728.70+145353.15 (Gaia DR3 1810475401285228928; $T_{\rm eff}=23{,}225$\,K,
$\log g=7.96$, $d=88$\,pc). TESS 20-s photometry from Sectors 55, 81 and 82 shows two modes at 334.85 and 283.55\,s and
their sum frequency. A search of public ZTF photometry at these two frequencies recovers both modes in $g$ and $r$
over 2018--2025 and finds nothing in a background star 3.9$''$ away, so the white dwarf is the pulsator. The ZTF
baseline gives periods of $334.84948\pm0.00008$\,s and $283.55430\pm0.00007$\,s. A possible third mode at 361.40\,s is
present in 2024 only. The star was flagged as variable by Gaia DR3 but listed as non-variable in a blind ZTF search;
we find no previous report of its pulsations.
\end{abstract}

\keywords{White dwarf stars (1799) --- DB stars (358) --- Stellar pulsations (1625) --- Asteroseismology (73)}

\section{Introduction}
Helium-atmosphere white dwarfs pulsate in non-radial gravity modes (DBV or V777~Her stars) in an instability strip
at roughly 21{,}000--30{,}000\,K \citep{corsico2022}. They are rare: \citet{vanderbosch2022} listed 28 previously
known DBVs and added 19 from SDSS spectroscopy, and \citet{corsico2022} added two from TESS. Each new DBV gives
periods that can constrain the mass, envelope and rotation of a white dwarf.

WD J203728.70+145353.15 (SDSS J203728.70+145353.0, TIC 456645161; $G=15.15$, parallax $11.40\pm0.04$\,mas) is
spectroscopically a DB with $T_{\rm eff}=23{,}225\pm1{,}087$\,K, $\log g=7.961\pm0.011$ and $M=0.582\,M_\odot$
\citep{kilic2020}, inside the DBV strip. Gaia DR3 flags it as variable, and the DR3 variability classifier
\citep{rimoldini2023} assigns it to the class of variable white dwarfs (score 0.84). Its G-band scatter is 15.7\,mmag,
4.6 times its expected error. We found it in a search of TESS data of white dwarfs within 100\,pc of the
\citet{gentilefusillo2021} catalogue.

\section{Data and analysis}
\textit{TESS.} We used the SPOC 20-s PDCSAP light curves of Sectors 55 (2022 August), 81 and 82 (2024 July--September).
A star of $G=14.85$ ($BP-RP=0.90$, parallax 1.18\,mas, i.e.\ a background star) lies 4.1$''$ away, and SPOC's crowding
correction assigns only 14--15\% of the aperture flux to the target (CROWDSAP 0.138--0.154); amplitudes below refer
to the white dwarf's flux after that correction. We kept QUALITY\,=\,0 points, removed a 0.25-d running median and
clipped at 5$\sigma$. False-alarm levels (0.1\%) come from the highest peak between 1{,}000 and 8{,}000\,$\mu$Hz in
shuffled light curves. Frequencies, amplitudes and phases were fitted by non-linear least squares with two free
frequencies and their sum; uncertainties are the formal ones, scaled by the residual scatter.

\textit{ZTF.} We retrieved the public ZTF PSF-photometry light curves \citep{bellm2019,masci2019} from IRSA within
1.5$''$ of the white dwarf and of the background star, which ZTF resolves as a separate source 3.9$''$ away. We
kept catflags\,=\,0 and exposures of $\le 60$\,s (1{,}097 $g$ and 1{,}790 $r$ epochs, 2018 April--2025 October;
300-s exposures would average a 335-s pulsation away). Magnitudes were converted to relative flux around the median of
each field and CCD quadrant, and times are heliocentric mid-exposure. Because ZTF sampling is sparse, we did not
search blindly: we computed the least-squares amplitude spectrum within $\pm1\,\mu$Hz of each TESS frequency and
compared its maximum with the same statistic at 300 control frequencies.

\section{Results}
The white dwarf shows two modes in all three TESS sectors (Figure~\ref{fig:1}a). Fitting Sectors 81+82 together gives
$f_1=2986.411\pm0.005\,\mu$Hz (334.8501\,s; $12.7\pm0.5$\,ppt) and $f_2=3526.657\pm0.006\,\mu$Hz (283.5546\,s;
$10.2\pm0.5$\,ppt). The sum frequency $f_1+f_2$ (153.54\,s) has $3.8\pm0.5$\,ppt, above the 0.1\% false-alarm amplitude
of 3.07\,ppt; such combination frequencies arise within a single pulsating star. The amplitudes of $f_1$ are
$13.1\pm0.7$, $14.4\pm0.7$ and $11.7\pm0.8$\,ppt in Sectors 55, 81 and 82, and its fitted frequency differs between
sectors by up to 0.07\,$\mu$Hz, several times the formal errors, which suggests small phase or frequency variations.
After removing $f_1$, $f_2$ and $f_1+f_2$, the highest residual peak in Sectors 81+82 is at $2767.04\,\mu$Hz (361.40\,s;
3.2\,ppt, 0.1\% false-alarm level 3.0\,ppt), present in Sectors 81 and 82 separately (3.2 and 3.3\,ppt) but not in
Sector 55 (1.2\,ppt). We regard it as a possible third mode.

ZTF recovers both TESS modes (Figure~\ref{fig:1}b,c). In $g$, $f_1$ and $f_2$ have $10.1\pm1.0$ and $9.0\pm1.0$\,ppt;
in $r$, $8.0\pm0.7$ and $5.8\pm0.7$\,ppt. All four peaks exceed the window maxima at all 300 control frequencies
(99.9th percentiles 4.5\,ppt in $g$ and 3.4\,ppt in $r$), and $f_1$ exceeds the 99th percentile in each of the eight
observing seasons in both bands. The $g$ and $r$ frequencies agree, and their weighted means,
$f_1=2986.4165\pm0.0007\,\mu$Hz and $f_2=3526.6614\pm0.0008\,\mu$Hz, agree with TESS within $1.1\sigma$. They correspond
to periods of $334.84948\pm0.00008$\,s and $283.55430\pm0.00007$\,s. The 361.40-s candidate is not detected in ZTF, whose
sensitivity ($\approx$4\,ppt) is not enough to test it. The background star shows 1.3--2.1\,ppt at both frequencies,
consistent with noise (false-alarm probabilities of 36--80\%). Together with the Gaia DR3 variability of the white
dwarf alone, this identifies the white dwarf as the pulsator.

\section{Discussion}
Three things probably kept this DBV unnoticed. In TESS the star contributes only $\sim$15\% of the light in its aperture.
The TESS-wide searches of \citet{romero2022,romero2025} were restricted to DA white dwarfs. And \citet{jestin2026},
who vetted Gaia DR3 variable white dwarfs with ZTF, list this star as not variable; their blind search detected
high-frequency pulsators with ZTF $g$ amplitudes of 0.07--0.11\,mag, several times the $\approx$0.01\,mag of this star
even allowing for differences in how amplitude is defined. A search at frequencies
already known from TESS reaches a few ppt with the same data. We tested it on 20 known TESS white-dwarf pulsators: for
periods below 400\,s it recovers 8 of 10 stars and every mode of 5.9\,ppt or more, with no detections in seven resolved
neighbours, while for periods of 400--1{,}000\,s it recovers only 4 of 10. It can therefore confirm short-period
candidates in crowded TESS apertures.

The ZTF amplitudes of the full 2018--2025 fit are lower than those of single seasons (11--18\,ppt), so some of the
difference from TESS may come from slow amplitude or phase changes; the cross-instrument amplitudes are also affected
by the TESS crowding correction and are not used here for mode identification. We searched SIMBAD, the ADS
(including full text), VizieR tables \citep{guidry2021,jestin2026} and the literature up to 2026 November 12 and found
no report of these pulsations. Fast ground-based photometry that resolves the 4$''$ pair would reveal more modes and
allow an asteroseismic analysis.

\begin{figure}[ht!]
\plotone{fig1.pdf}
\caption{(a) TESS amplitude spectrum of WD J203728.70+145353.15 (Sectors 81+82, 20-s cadence) with the 0.1\%
false-alarm level; the possible third mode is marked ``361.4\,s?''. (b, c) ZTF least-squares amplitude spectra around
the two TESS frequencies for the white dwarf in $g$ (green) and $r$ (red), and for the background star 3.9$''$ away
in $g$ (dotted grey).\label{fig:1}}
\end{figure}

\begin{acknowledgments}
This paper includes data collected by the TESS mission, obtained from MAST; funding for the TESS mission is provided by
the NASA Science Mission Directorate. It is based on observations obtained with the Samuel Oschin Telescope 48-inch
and the 60-inch Telescope at the Palomar Observatory as part of the Zwicky Transient Facility project, obtained
from IRSA. It has made use of data from the European Space Agency (ESA) mission Gaia, processed by the Gaia Data
Processing and Analysis Consortium, and of the SIMBAD and VizieR services and NASA's Astrophysics Data System.
We made use of AI-assisted tools (Anthropic's Claude) to speed up data retrieval, code development and the
preparation of this note; all analyses were verified by the author, who is responsible for the content.
\end{acknowledgments}

\facilities{TESS, PO:1.2m (ZTF), Gaia}
\software{astropy, astroquery, lightkurve, matplotlib, numpy, scipy}

\bibliography{note}{}
\bibliographystyle{aasjournal}

\end{document}
